% ============================================================
\section{Giới thiệu}
% ============================================================
Sự bùng nổ của dữ liệu văn bản số tại Việt Nam -- từ báo điện
tử, thương mại điện tử, mạng xã hội đến kho lưu trữ học thuật -- khiến việc tổ
chức và phân loại thủ công ngày càng bất khả thi, đặt ra yêu cầu về các hệ thống
tự động hóa. Trong số đó, \emph{phân loại văn bản đa
nhãn zero-shot} (zero-shot multi-label text classification) trong môi trường thế
giới mở nổi lên như một bài toán then chốt: hệ thống tự động xác định không gian
nhãn (label space) và gán nhiều danh mục cho mỗi văn bản mà không cần bất kỳ dữ
liệu gán nhãn theo tác vụ hay tập nhãn định trước nào~\cite{li2024open}. Ở đây,
``thế giới mở'' (open-world) chỉ một môi trường động, nơi không gian nhãn không
được định trước mà được khám phá và mở rộng dần khi các batch dữ liệu mới đến
theo thời gian.

Một tình huống thực tế điển hình đến từ các công cụ quản lý tri thức cá nhân như
NotebookLM: người dùng tải lên hàng loạt tài liệu nhưng thiếu cơ chế tự động
phân cụm và phân cấp, buộc phải quản lý thủ công bằng thư mục hoặc đánh dấu chọn,
gây bất tiện khi số lượng tài liệu lớn.

Các hướng tiếp cận hiện có cho bài toán
này có thể chia thành ba nhóm~\cite{li2024open}. \emph{Thứ nhất}, học có giám sát
truyền thống cho độ chính xác cao nhưng cần tập nhãn đầy đủ và dữ liệu gán nhãn
-- trái với giả định không nhãn. \emph{Thứ hai}, dùng LLM để sinh nhãn trực tiếp
trên văn bản thô: linh hoạt nhưng chi phí gọi API rất lớn và dễ sinh nhãn nhiễu
do hạn chế cửa sổ ngữ cảnh. \emph{Thứ ba}, các phương pháp thống kê (mô hình chủ
đề, trích xuất từ khóa): nhanh, nhẹ, chi phí thấp nhưng chỉ trả về từ/cụm từ rời
rạc, thiếu tính khái quát để làm tên nhãn.

Từ đó nổi lên một khoảng trống rõ rệt: cần một giải pháp dung hòa được
\emph{chất lượng}, \emph{chi phí} và \emph{tốc độ}, đặc biệt cho tiếng Việt --
ngôn ngữ đơn lập, giàu từ đồng nghĩa và chưa được tối ưu tốt bởi các mô hình đa
ngữ. Bài báo này đề xuất \textbf{StatLM} (\emph{Statistical Encoder -- Language
Model Decoder}), một kiến trúc pipeline tách biệt vai trò giữa thành phần thống
kê và LLM. Ý tưởng cốt lõi: một \emph{encoder thống kê} đảm nhiệm giảm nhiễu và
cô đặc thông tin cốt lõi từ văn bản thô, trong khi \emph{decoder LLM} chỉ được
huy động tập trung một lần ở khâu đặt tên nhãn cuối cùng. Sự phân tách này cho
phép StatLM tận dụng đồng thời tốc độ của kỹ thuật thống kê và sức mạnh khái quát
hóa ngữ nghĩa của LLM.

Bài báo đặt ra ba câu hỏi nghiên cứu. \textbf{RQ1:} Những hạn chế nào tồn tại
trong các phương pháp phân loại đa nhãn zero-shot hiện nay, và kiến trúc
encoder--decoder khắc phục chúng ra sao? \textbf{RQ2:} Việc thay hàm tương đồng
từ vựng của TextRank bằng sentence embedding và tùy chỉnh KeyBERT cho tiếng Việt
cải thiện chất lượng tóm tắt và trích xuất keyphrase đến mức nào? \textbf{RQ3:}
Với việc chỉ dùng LLM ở decoder, StatLM cân bằng giữa chất lượng không gian nhãn,
chi phí và tốc độ ra sao so với hướng dùng LLM toàn phần?

\textbf{Đóng góp.} (i) Hệ thống hóa hạn chế của ba nhóm phương pháp và định vị
khoảng trống nghiên cứu, từ đó đề xuất kiến trúc encoder--decoder. (ii) Thiết kế
và triển khai pipeline StatLM hoàn chỉnh cho tiếng Việt, trong đó từng tầng được
lựa chọn công nghệ một cách có nguyên lý (thay hàm tương đồng của TextRank, cấu
hình KeyBERT bất đối xứng). (iii) Xây dựng và công bố công khai bộ dữ liệu
VietnamNet, cùng đánh giá thực nghiệm định lượng và định tính chứng minh hiệu quả
của giải pháp.

